% Complete Section 3.2, source PDF 32-33 / printed 22-23.
% block: closed-operators-heading page:32 type:section
\subsection{\BT{Closed operators}{闭算子}}\label{sec:closed-operators}

% block: closed-operator-definition pages:32,33 type:ordinary-paragraph
\BT{Many important properties of bounded linear operators extend to a much larger class of unbounded operators, called \emph{closed}. Banach's closed graph theorem states than an everywhere defined linear operator between Hilbert (or Banach) spaces is bounded if and only if its graph is closed. This inspires the definition of a \emph{closed operator}, which is any unbounded operator whose graph is closed in $X\times Y$. In other words, an unbounded operator $T$ from $X$ to $Y$ is a closed linear operator if and only if the following condition holds:}{有界线性算子的许多重要性质可推广到一类大得多的无界算子, 即\emph{闭算子}. Banach 闭图定理指出, Hilbert (或 Banach) 空间之间处处有定义的线性算子有界, 当且仅当其图是闭的. 由此引出\emph{闭算子}的定义, 即其图在 $X\times Y$ 中闭的无界算子. 换言之, 从 $X$ 到 $Y$ 的无界算子 $T$ 是闭线性算子, 当且仅当下述条件成立:}

% block: closed-operator-sequence-criterion page:33 type:ordinary-paragraph
\begin{quote}
\BT{Whenever $v_1,v_2,\ldots\in D(T)$ satisfy $v_n\xrightarrow{X}x$ and $Tv_n\xrightarrow{Y}y$ for some $x\in X$ and $y\in Y$, then $x\in D(T)$ and $Tx=y$.}{只要 $v_1,v_2,\ldots\in D(T)$ 对某些 $x\in X$ 和 $y\in Y$ 满足 $v_n\xrightarrow{X}x$ 及 $Tv_n\xrightarrow{Y}y$, 就有 $x\in D(T)$ 且 $Tx=y$.}
\end{quote}

% block: graph-complete-criterion-intro page:33 type:ordinary-paragraph
\BT{The following proposition gives another, equivalent definition of a closed operator.}{下面的命题给出闭算子的另一个等价定义.}

% block: graph-complete-criterion page:33 type:proposition
\Needspace{6\baselineskip}
\begin{Proposition}\label{prop:closed-graph-domain-complete}
\BT{An unbounded operator $T$ is closed if and only if its domain, endowed with the graph norm, is a Hilbert space.}{无界算子 $T$ 是闭的, 当且仅当其定义域配备图范数后成为 Hilbert 空间.}
\end{Proposition}

% block: graph-complete-proof-forward page:33 type:proof
\begin{Proof}
\BT{We must show that $D(T)$ is complete if and only if $T$ is closed. Let $\{v_n\}_{n\in\mathbb N}$ be a sequence of elements of $D(T)$. To say the sequence is Cauchy with respect to the graph norm means that it is Cauchy with respect to the $X$ norm and $\{Tv_n\}$ is Cauchy with respect to the $Y$ norm. Since $X$ and $Y$ are complete, this implies that $v_n\xrightarrow{X}x$ and $Tv_n\xrightarrow{Y}y$ for some $x\in X$, $y\in Y$. If $T$ is closed, then $x\in D(T)$ and $Tx=y$, so $v_n\xrightarrow{D(T)}x$. Thus $D(T)$ is complete.}{须证明 $D(T)$ 完备当且仅当 $T$ 是闭的. 设 $\{v_n\}_{n\in\mathbb N}$ 是 $D(T)$ 中的一个元素序列. 该序列关于图范数为 Cauchy 序列, 意味着它关于 $X$ 范数为 Cauchy 序列, 且 $\{Tv_n\}$ 关于 $Y$ 范数为 Cauchy 序列. 由于 $X$ 和 $Y$ 完备, 这意味着对某些 $x\in X$, $y\in Y$, 有 $v_n\xrightarrow{X}x$ 及 $Tv_n\xrightarrow{Y}y$. 若 $T$ 是闭的, 则 $x\in D(T)$ 且 $Tx=y$, 从而 $v_n\xrightarrow{D(T)}x$. 因此 $D(T)$ 完备.}

% block: graph-complete-proof-reverse page:33 type:proof
\BT{Conversely, suppose $D(T)$ is complete with respect to the graph norm and that $v_n\xrightarrow{X}x$ and $Tv_n\xrightarrow{Y}y$. Then the sequence $v_n$ is Cauchy with respect to the graph norm on $D(T)$, so $v_n\xrightarrow{D(T)}v$ for some $v\in D(T)$, i.e., $v_n\xrightarrow{X}v$ and $Tv_n\xrightarrow{Y}Tv$. Thus $x=v\in D(T)$ and $y=Tv=Tx$, so $T$ is closed.}{反之, 设 $D(T)$ 关于图范数完备, 且 $v_n\xrightarrow{X}x$ 及 $Tv_n\xrightarrow{Y}y$. 则序列 $v_n$ 关于 $D(T)$ 上的图范数为 Cauchy 序列, 从而对某个 $v\in D(T)$ 有 $v_n\xrightarrow{D(T)}v$, 即 $v_n\xrightarrow{X}v$ 且 $Tv_n\xrightarrow{Y}Tv$. 因此 $x=v\in D(T)$, $y=Tv=Tx$, 故 $T$ 是闭的.}
\end{Proof}
